% country: BLR
Let $n$ be an integer greater than 2. A positive integer is said to be \emph{attainable} if it is 1 or can be obtained from 1 by a sequence of operations with the following properties:
\begin{itemize}
\item[(i)] The first operation is either addition or multiplication.
\item[(ii)] Thereafter, additions and multiplications are used alternately.
\item[(iii)] In each addition one can choose independently whether to add 2 or $n$.
\item[(iv)] In each multiplication, one can choose independently whether to multiply by 2 or by $n$.
\end{itemize}
A positive integer that cannot be so obtained is said to be \emph{unattainable}.
\begin{itemize}
\item[(a)] Prove that if $n \geq 9$, there are infinitely many unattainable positive integers.
\item[(b)] Prove that if $n = 3$, all positive integers except 7 are attainable.
\end{itemize}
